% !TEX root = ../proyect.tex

\chapter{Resultados y discusión}\label{cap:resultados}

\section{Resultados de la caracterización inicial}\label{sec:resultados-caracterizacion}
% TODO: redactar este apartado.

\section{Resultados del filtrado lineal}\label{sec:resultados-filtrado}
% TODO: redactar este apartado.

\section{Evidencia de dependencia no lineal}\label{sec:evidencia-no-lineal}
% TODO: redactar este apartado.

\section{Resultados de la reconstrucción del espacio de estados}\label{sec:resultados-reconstruccion}
% TODO: redactar este apartado.

\section{Resultados de cuantificación dinámica}\label{sec:resultados-cuantificacion}
% TODO: redactar este apartado.

\section{Resultados de predicción local}\label{sec:resultados-prediccion}
% TODO: redactar este apartado.

\section{Comparación con modelos de referencia}\label{sec:comparacion-referencia}
% TODO: redactar este apartado.

\section{Comparación entre configuraciones alternativas}\label{sec:comparacion-configuraciones}
% TODO: redactar este apartado.

\section{Comparación final con HAR-logRV}\label{sec:resultados-har}
La comparación final incorpora HAR-logRV como modelo práctico de memoria multiescala. En la muestra test comparable de 5000 puntos, HAR-logRV obtiene un RMSE de 0,8608, frente a 0,8641 de AR(49) y 0,8765 del kNN local con \(\tau=137\), \(m=5\) y \(k=200\). La mejora frente a AR(49) es pequeña, aproximadamente un 0,4\%, por lo que no debe presentarse como una ruptura fuerte respecto al benchmark lineal. Su interés principal está en que combina buen rendimiento, bajo coste computacional, interpretabilidad y facilidad de exportación al MVP.

Este resultado también matiza la lectura de las fases de reconstrucción: el kNN confirma cierta utilidad predictiva de la estructura local, pero en BTC la persistencia multiescala de la volatilidad queda mejor resumida por un modelo HAR simple. Por tanto, HAR-logRV se interpreta como cierre práctico del estudio, no como evidencia adicional de caos determinista.

\section{Discusión sobre la existencia de caos determinista}\label{sec:discusion-caos}
% TODO: redactar este apartado.

\section{Interpretación global de los resultados}\label{sec:interpretacion-global}
% TODO: redactar este apartado.
